\sectionkicker{Source-backed technical notes}
\subsection*{ModelからActionへ — Operator・Deep Research・Responses API・Claude Code・Codex・Jules: Technical Notes}
本欄は記事本文で圧縮した一次資料上の情報を、比較・再検証しやすい形で整理したものである。時系列、確認済みの主張、留意点、一次資料URLを掲載する。
\medskip
\subsection*{Theme at a glance}
\addcontentsline{toc}{subsection}{Theme at a glance}
\begin{center}
\footnotesize
\begin{tabularx}{\linewidth}{@{}p{0.20\linewidth}p{0.10\linewidth}p{0.14\linewidth}X@{}}
\toprule
資料 & 位置づけ & 種別 & 時系列 \\
\midrule
OpenAI Operator research preview & 主要資料 & - & 2025-01-23   \\
OpenAI deep research & 主要資料 & - & 2025-02-02   \\
Responses API, built-in tools, and Agents SDK & 主要資料 & - & 2025-03-11   \\
OpenAI Codex cloud software engineering agent & 主要資料 & - & 2025-05-16   \\
Jules coding agent public beta & 補足資料 & - & 2025-05-20   \\
Responses API remote MCP and tool expansion & 補足資料 & - & 2025-05-21   \\
Anthropic web search API & 補足資料 & - & 2025-05-07   \\
\bottomrule
\end{tabularx}
\end{center}
\normalsize

\begin{technicalnote}{OpenAI Operator research preview}{主要資料}
\begingroup
% reader-facing Technical Notes break policy
\widowpenalty=10000
\clubpenalty=10000
\displaywidowpenalty=10000
\begin{tabularx}{\linewidth}{@{}>{\bfseries}p{0.22\linewidth}X@{}}
組織 & OpenAI \\
種別 & - \\
時系列 & 2025-01-23 \\
\end{tabularx}
\smallskip
{\bfseries 一次資料から整理したtechnical points}
\begin{itemize}[leftmargin=1.5em,itemsep=0.35em]
\item \textbf{一次情報で確認できる事実}: 2025年1月23日、OpenAIは米国のProユーザー向けにブラウザ操作Agent「Operator」のresearch previewを開始した。
\end{itemize}
\begin{samepage}
{\bfseries 一次資料}
\begingroup\sloppy
\Urlmuskip=0mu plus 2mu\relax
\begin{itemize}[leftmargin=1.5em,itemsep=0.25em]
\item {\scriptsize\url{https://openai.com/index/introducing-operator/}}
\end{itemize}
\endgroup
\end{samepage}
\endgroup
\end{technicalnote}
\medskip

\begin{technicalnote}{OpenAI deep research}{主要資料}
\begingroup
% reader-facing Technical Notes break policy
\widowpenalty=10000
\clubpenalty=10000
\displaywidowpenalty=10000
\begin{tabularx}{\linewidth}{@{}>{\bfseries}p{0.22\linewidth}X@{}}
組織 & OpenAI \\
種別 & - \\
時系列 & 2025-02-02 \\
\end{tabularx}
\smallskip
{\bfseries 一次資料から整理したtechnical points}
\begin{itemize}[leftmargin=1.5em,itemsep=0.35em]
\item \textbf{一次情報で確認できる事実}: 2025年2月2日、OpenAIはChatGPTで複数段階のWeb調査を行うDeep Researchを公開した。
\end{itemize}
\begin{samepage}
{\bfseries 一次資料}
\begingroup\sloppy
\Urlmuskip=0mu plus 2mu\relax
\begin{itemize}[leftmargin=1.5em,itemsep=0.25em]
\item {\scriptsize\url{https://openai.com/index/introducing-deep-research/}}
\end{itemize}
\endgroup
\end{samepage}
\endgroup
\end{technicalnote}
\medskip

\begin{technicalnote}{Responses API, built-in tools, and Agents SDK}{主要資料}
\begingroup
% reader-facing Technical Notes break policy
\widowpenalty=10000
\clubpenalty=10000
\displaywidowpenalty=10000
\begin{tabularx}{\linewidth}{@{}>{\bfseries}p{0.22\linewidth}X@{}}
組織 & OpenAI \\
種別 & - \\
時系列 & 2025-03-11 \\
\end{tabularx}
\smallskip
{\bfseries 一次資料から整理したtechnical points}
\begin{itemize}[leftmargin=1.5em,itemsep=0.35em]
\item \textbf{一次情報で確認できる事実}: 2025年3月11日、OpenAIはResponses API、Web/File Search、Computer Use、Agents SDK、TracingなどのAgent構築機能を発表した。
\end{itemize}
\begin{samepage}
{\bfseries 一次資料}
\begingroup\sloppy
\Urlmuskip=0mu plus 2mu\relax
\begin{itemize}[leftmargin=1.5em,itemsep=0.25em]
\item {\scriptsize\url{https://openai.com/index/new-tools-for-building-agents/}}
\end{itemize}
\endgroup
\end{samepage}
\endgroup
\end{technicalnote}
\medskip

\begin{technicalnote}{OpenAI Codex cloud software engineering agent}{主要資料}
\begingroup
% reader-facing Technical Notes break policy
\widowpenalty=10000
\clubpenalty=10000
\displaywidowpenalty=10000
\begin{tabularx}{\linewidth}{@{}>{\bfseries}p{0.22\linewidth}X@{}}
組織 & OpenAI \\
種別 & - \\
時系列 & 2025-05-16 \\
\end{tabularx}
\smallskip
{\bfseries 一次資料から整理したtechnical points}
\begin{itemize}[leftmargin=1.5em,itemsep=0.35em]
\item \textbf{一次情報で確認できる事実}: 2025年5月16日、OpenAIは並列sandboxでタスクを実行するcloud software-engineering Agent Codexをresearch previewとして公開した。
\end{itemize}
\begin{samepage}
{\bfseries 一次資料}
\begingroup\sloppy
\Urlmuskip=0mu plus 2mu\relax
\begin{itemize}[leftmargin=1.5em,itemsep=0.25em]
\item {\scriptsize\url{https://openai.com/index/introducing-codex/}}
\end{itemize}
\endgroup
\end{samepage}
\endgroup
\end{technicalnote}
\medskip

\begin{technicalnote}{Jules coding agent public beta}{補足資料}
\begingroup
% reader-facing Technical Notes break policy
\widowpenalty=10000
\clubpenalty=10000
\displaywidowpenalty=10000
\begin{tabularx}{\linewidth}{@{}>{\bfseries}p{0.22\linewidth}X@{}}
組織 & Google \\
種別 & - \\
時系列 & 2025-05-20 \\
\end{tabularx}
\smallskip
{\bfseries 一次資料から整理したtechnical points}
\begin{itemize}[leftmargin=1.5em,itemsep=0.35em]
\item \textbf{一次情報で確認できる事実}: 2025年5月20日、GoogleはGitHub repositoryをcloud VM上で扱う非同期coding Agent Julesをpublic betaで公開した。
\end{itemize}
\begin{samepage}
{\bfseries 一次資料}
\begingroup\sloppy
\Urlmuskip=0mu plus 2mu\relax
\begin{itemize}[leftmargin=1.5em,itemsep=0.25em]
\item {\scriptsize\url{https://developers.googleblog.com/google-io-2025-developer-keynote-recap/}}
\end{itemize}
\endgroup
\end{samepage}
\endgroup
\end{technicalnote}
\medskip

\begin{technicalnote}{Responses API remote MCP and tool expansion}{補足資料}
\begingroup
% reader-facing Technical Notes break policy
\widowpenalty=10000
\clubpenalty=10000
\displaywidowpenalty=10000
\begin{tabularx}{\linewidth}{@{}>{\bfseries}p{0.22\linewidth}X@{}}
組織 & OpenAI \\
種別 & - \\
時系列 & 2025-05-21 \\
\end{tabularx}
\smallskip
{\bfseries 一次資料から整理したtechnical points}
\begin{itemize}[leftmargin=1.5em,itemsep=0.35em]
\item \textbf{一次情報で確認できる事実}: 2025年5月21日、OpenAIはResponses APIにremote MCP、image generation、Code Interpreter、background mode、reasoning summariesを追加した。
\end{itemize}
\begin{samepage}
{\bfseries 一次資料}
\begingroup\sloppy
\Urlmuskip=0mu plus 2mu\relax
\begin{itemize}[leftmargin=1.5em,itemsep=0.25em]
\item {\scriptsize\url{https://openai.com/index/new-tools-and-features-in-the-responses-api/}}
\end{itemize}
\endgroup
\end{samepage}
\endgroup
\end{technicalnote}
\medskip

\begin{technicalnote}{Anthropic web search API}{補足資料}
\begingroup
% reader-facing Technical Notes break policy
\widowpenalty=10000
\clubpenalty=10000
\displaywidowpenalty=10000
\begin{tabularx}{\linewidth}{@{}>{\bfseries}p{0.22\linewidth}X@{}}
組織 & Anthropic \\
種別 & - \\
時系列 & 2025-05-07 \\
\end{tabularx}
\smallskip
{\bfseries 一次資料から整理したtechnical points}
\begin{itemize}[leftmargin=1.5em,itemsep=0.35em]
\item \textbf{一次情報で確認できる事実}: 2025年5月7日、AnthropicはMessages APIにWeb Search toolを追加した。
\end{itemize}
\begin{samepage}
{\bfseries 一次資料}
\begingroup\sloppy
\Urlmuskip=0mu plus 2mu\relax
\begin{itemize}[leftmargin=1.5em,itemsep=0.25em]
\item {\scriptsize\url{https://www.anthropic.com/news/web-search-api}}
\end{itemize}
\endgroup
\end{samepage}
\endgroup
\end{technicalnote}
\medskip
